% !TeX root = ActuarialFormSheet_MBFA-es.tex
% !TeX spellcheck = es_ES

%\subsection*{Lois utilisées dans la théorie des extrêmes}


\begin{f}[Ley de Pareto]
	Sea la variable aleatoria \(X\) que sigue una ley de Pareto de parámetros \((x_{\mathrm{m}},k)\), siendo \(k\) el índice de Pareto:
	\[
	\mathbb{P}(X>x)=\left(\frac{x}{x_{\mathrm{m}}}\right)^{-k}\ \mbox{con}\ x \geq x_{\mathrm{m}} 
	\]
	
	\[
	f_{k,x_\mathrm{m}}(x) = k\,\frac{x_\mathrm{m}^k}{x^{k+1}}\ \mbox{para}\ x \ge x_\mathrm{m}
	\]
	
	\textbf{Ley de Pareto generalizada} (LPG) con 3 parámetros \(\mu\), \(\sigma\) y \(\xi\).
	
	
	\[
	F_{\xi,\mu,\sigma}(x) = \begin{cases} 1 - \left(1+ \frac{\xi(x-\mu)}{\sigma}\right)^{-1/\xi} & \text{para }\xi \neq 0, \\ 1 - \exp \left(-\frac{x-\mu}{\sigma}\right) & \text{para }\xi = 0. \end{cases} 
	\]
	para \(x \geq \mu\) cuando \(\xi \geq 0\) y \( \mu \leq x \leq \mu - \sigma /\xi\) cuando \( \xi < 0\), donde \(\mu\in\mathbb{R}\) es la localización, \(\sigma>0\) la escala y  \(\xi\in\mathbb{R}\) la forma. 
	Cabe señalar que algunas referencias dan el parámetro de forma como \(\kappa = - \xi\).

\[
f_{\xi,\mu,\sigma}(x) = \frac{1}{\sigma}\left(1 + \frac{\xi (x-\mu)}{\sigma}\right)^{\left(-\frac{1}{\xi} - 1\right)}
=
 \frac{\sigma^{\frac{1}{\xi}}}{\left(\sigma + \xi (x-\mu)\right)^{\frac{1}{\xi}+1}}
\]
%à nouveau, pour  \(x \geq \mu\), et \(x \leq \mu - \sigma /\xi\) où \(\xi < 0\). 

\medskip

\begin{center}
		
\tikzset{
	declare function={
		dParetoGen(\x,\mu,\sigma,\xi) = 
		(\x>\mu)*(1/\sigma*(1 + \xi* (\x-\mu)/\sigma)^(-1/\xi - 1));}
}  
\begin{tikzpicture}[mark=none,samples=100,smooth,domain=0:15]
	\begin{axis}[ ytick=\empty,xtick pos=left,
		axis y line=middle,
		axis x line=bottom,
		height=.5*\linewidth,width=0.95*\linewidth, ticklabel style ={font=\footnotesize}, legend pos=north east,
		legend style={font=\tiny}]
		\addplot[BleuProfondIRA] 	({\x},{dParetoGen({\x},1,1,1)}) ;
		\addplot[FushiaIRA] 	({\x},{dParetoGen({\x},1,1,5)}) ;
		\addplot[VertIRA]  	({\x},{dParetoGen({\x},1,1,20)}) ;
		\addplot[OrangeProfondIRA] 	({\x},{dParetoGen({\x},3,2,5)}) ;
		%
		\legend{\(\mu\) =  1 \(\sigma=1\) \(\xi = 1\),\(\mu\) =  1  \(\sigma=2\) \(\xi =5\),\(\mu\) =  1  \(\sigma=3\) \(\xi =20\),\(\mu\) = 3  \(\sigma=2\) \(\xi = 5\)};
	\end{axis} 
\end{tikzpicture}

\end{center}


%\begin{center}
%	{\shorthandoff{!;}	
%	\tikzset{
%	declare function={
%		pParetoGen(\x,\mu,\sigma,\xi) =  (\x > \mu)*( 1-(1+ (\xi*(\x-\mu))/\sigma)^(-1/\xi));}
%}  
%\begin{tikzpicture}[mark=none,samples=100,smooth,domain=0:15]
%	\begin{axis}[ ytick=\empty,xtick pos=left,
%		axis y line=middle,
%		axis x line=bottom,
%		height=.5*\linewidth,width=0.95*\linewidth, ticklabel style ={font=\footnotesize}, legend pos=north east,
%		legend style={font=\tiny}]
%		\addplot[BleuProfondIRA] 
%		({\x},{pParetoGen({\x},1,1,1)}) ;
%		\addplot[FushiaIRA] 
%		({\x},{pParetoGen({\x},1,1,5)}) ;
%		\addplot[VertIRA] 
%		({\x},{pParetoGen({\x},1,1,20)}) ;
%		\addplot[OrangeProfondIRA] 
%		({\x},{pParetoGen({\x},3,2,5)}) ;
%		%
%		\legend{\(\mu\) =  1 \(\sigma=1\) \(\xi = 1\),\(\mu\) =  1  \(\sigma=2\) \(\xi =5\),\(\mu\) =  1  \(\sigma=3\) \(\xi =20\),\(\mu\) = 3  \(\sigma=2\) \(\xi = 5\)};
%	\end{axis}	 
%\end{tikzpicture}
%}
%\end{center}

\end{f}
\hrule

\begin{f}[Distribución de valores extremos generalizada]
	La función de distribución de la distribución de valores extremos generalizada es:
	\[
	F_{\mu,\sigma,\xi}(x) = \exp\left\{-\left[1+\xi\left(\frac{x-\mu}{\sigma}\right)\right]^{-1/\xi}\right\}
	\]
	para \(1+\xi(x-\mu)/\sigma>0\), donde \(\mu\in\mathbb{R}\) es la localización, \(\sigma>0\) la escala y  \(\xi\in\mathbb{R}\) la forma. Para \(\xi = 0\), la expresión viene definida por su límite en 0.

\begin{align*}
	f_{\mu,\sigma,\xi}(x) =& \frac{1}{\sigma}\left[1+\xi\left(\frac{x-\mu}{\sigma}\right)\right]^{(-1/\xi)-1}\\ &\times\exp\left\{-\left[1+\xi\left(\frac{x-\mu}{\sigma}\right)\right]^{-1/\xi}\right\}
\end{align*}

\[
f(x;\mu ,\sigma ,0)={\frac {1}{\sigma }}\exp \left(-{\frac {x-\mu }{\sigma }}\right)\exp \left[-\exp \left(-{\frac {x-\mu }{\sigma }}\right)\right]
\]

\end{f}
\hrule



\begin{f}[Ley de Gumbel]
	
	La función de distribución de la \textbf{ley de Gumbel} es:
	\[
	F_{\mu,\sigma}(x) = e^{-e^{(\mu-x)/\sigma}}.\,
	\]
	Para \(\mu = 0 \) y \(\sigma = 1\), se obtiene la ley de Gumbel estándar.
	La ley de Gumbel es un caso particular de la ley GEV (con \(\xi=0\)).
	
	Su densidad:
\[
f_{\mu,\sigma}(x)=\frac{1}{\sigma}e^{\left(\frac{x-\mu}{\sigma}-e^{-(x-\mu)/\sigma}\right)}
\]


\begin{center}
\begin{tikzpicture}[mark=none,samples=100,smooth,domain=-5:15]
\begin{axis}[ ytick=\empty,xtick pos=left,
	axis y line=middle,
	axis x line=bottom,
	height=.5*\linewidth,width=0.95*\linewidth, ticklabel style ={font=\footnotesize}, legend pos=north east,
	legend style={font=\tiny}]
	\addplot[BleuProfondIRA] 
	({\x},{dGumbel({\x},0.5,1)}) ;
	\addplot[FushiaIRA] 
	plot ({\x},{dGumbel({\x},1,2)}) ;
	\addplot[VertIRA] 
	plot ({\x},{dGumbel({\x},1.5,3)}) ;
	\addplot[OrangeProfondIRA] 
	plot ({\x},{dGumbel({\x},3,4)}) ;
	%
	\legend{\(\mu\) =  0.5 \(\gamma = 1\),\(\mu\) =  1 \(\gamma =2\),\(\mu\) =  1.5 \(\gamma =3\),\(\mu\) =  3 \(\gamma = 4\)};
\end{axis}	 
\end{tikzpicture}
\end{center}


%\begin{center}
%\begin{tikzpicture}[mark=none,samples=100,smooth,domain=-5:15]
%\begin{axis}[ ytick=\empty,xtick pos=left,
%	axis y line=middle,
%	axis x line=bottom,
%	height=.5*\linewidth,width=0.95*\linewidth, ticklabel style ={font=\footnotesize}, legend pos=north west,
%	legend style={font=\tiny}]
%	\addplot[BleuProfondIRA] 
%	({\x},{pGumbel({\x},0.5,1)}) ;
%	\addplot[FushiaIRA] 
%	plot ({\x},{pGumbel({\x},1,2)}) ;
%	\addplot[VertIRA] 
%	plot ({\x},{pGumbel({\x},1.5,3)}) ;
%	\addplot[OrangeProfondIRA] 
%	plot ({\x},{pGumbel({\x},3,4)}) ;
%	%
%	\legend{\(\mu\) =  0.5 \(\gamma = 1\),\(\mu\) =  1 \(\gamma =2\),\(\mu\) =  1.5 \(\gamma =3\),\(\mu\) =  3 \(\gamma = 4\)};
%\end{axis}	 
%\end{tikzpicture}
%\end{center}

%\subsubsubsection{Loi de Weibull}\label{label}
\end{f}
\hrule

\begin{f}[Ley de Weibull]
	
	
	La \textbf{ley de Weibull} tiene como función de distribución la siguiente:
	\[
	F_{\alpha,\mu, \sigma}(x) = 1- e^{-((x-\mu)/\sigma)^\alpha}\,
	\]
	donde \(x >\mu\).
	Su densidad de probabilidad es:
	\[
	f_{\alpha,\mu,\sigma}(x) = (\alpha/\sigma) ((x-\mu)/\sigma)^{(\alpha-1)} e^{-((x-\mu)/\sigma)^\alpha}\,
	\]
	donde  \(\mu\in\mathbb{R}\) es la localización, \(\sigma>0\) la escala y  \(\alpha=-1/\xi>0 \) la forma.
	
	La distribución de Weibull se utiliza a menudo en el ámbito del análisis de la vida útil. Es un caso particular de la distribución GEV cuando \(\xi<0\).
	
	Si la tasa de fallos disminuye con el tiempo, entonces \(\alpha<1\). Si la tasa de fallos es constante a lo largo del tiempo, entonces \(\alpha=1\). Si la tasa de fallos aumenta con el tiempo, entonces \(\alpha>1\). Comprender la tasa de fallos puede proporcionar una indicación sobre la causa de los fallos.
	

\begin{center}
	\tikzset{
	declare function={
		dWeibull(\x,\sigma,\alpha) = (\alpha/\sigma)*(\x/\sigma)^(\alpha-1)*exp(-(\x/\sigma)^\alpha);}
}  
\begin{tikzpicture}[mark=none,samples=100,smooth,domain=0:5]
	\begin{axis}[ ytick=\empty,xtick pos=left,
		axis y line=middle,
		axis x line=bottom,
		height=.5*\linewidth,width=0.95*\linewidth, ticklabel style ={font=\footnotesize}, legend pos=north east,
		legend style={font=\tiny}]
		\addplot[BleuProfondIRA] 
		({\x},{dWeibull({\x},0.5,1)}) ;
		\addplot[FushiaIRA] 
		plot ({\x},{dWeibull({\x},1,2)}) ;
		\addplot[VertIRA] 
		plot ({\x},{dWeibull({\x},1.5,3)}) ;
		\addplot[OrangeProfondIRA] 
		plot ({\x},{dWeibull({\x},3,4)}) ;
		%
		\legend{\(\mu\) =  0.5 \(\gamma = 1\) (\(\alpha=-1/\xi>0\)),\(\mu\) =  1 \(\gamma =2\),\(\mu\) =  1.5 \(\gamma =3\),\(\mu\) =  3 \(\gamma = 4\)};
	\end{axis}	 
\end{tikzpicture}

%	{\shorthandoff{!;}	
%		\tikzset{
%	declare function={
%		pWeibull(\x,\sigma,\alpha) = 1- exp(-(\x/\sigma)^\alpha);}
%}  
%\begin{tikzpicture}[mark=none,samples=100,smooth,domain=0:5]
%	\begin{axis}[ ytick=\empty,xtick pos=left,
%		axis y line=middle,
%		axis x line=bottom,
%		height=.5*\linewidth,width=0.95*\linewidth, ticklabel style ={font=\footnotesize}, legend pos= south east,
%		legend style={font=\tiny}]
%		\addplot[BleuProfondIRA] 
%		({\x},{pWeibull({\x},0.5,1)}) ;
%		\addplot[FushiaIRA] 
%		plot ({\x},{pWeibull({\x},1,2)}) ;
%		\addplot[VertIRA] 
%		plot ({\x},{pWeibull({\x},1.5,3)}) ;
%		\addplot[OrangeProfondIRA] 
%		plot ({\x},{pWeibull({\x},3,4)}) ;
%		%
%		\legend{\(\mu\) =  0.5 \(\gamma = 1\) (\(\alpha=-1/\xi>0\)),\(\mu\) =  1 \(\gamma =2\),\(\mu\) =  1.5 \(\gamma =3\),\(\mu\) =  3 \(\gamma = 4\)};
%	\end{axis}	 
%\end{tikzpicture}
%}
\end{center}
\end{f}


\begin{f}[Ley de Fréchet]
	La función de distribución de la \textbf{ley de Fréchet}\index{Definición: ley de Fréchet} viene dada por 
	\[
	F_{\alpha,\mu,\sigma}(x)=\mathbb{P}(X \le x)=\begin{cases} e^{-\left(\frac{x-\mu}{\sigma}\right)^{-\alpha}} & \text{ si } x>\mu \\ 0 &\text{ en caso contrario.}\end{cases} 
	\]
	donde  \(\mu\in\mathbb{R}\) es la localización, \(\sigma>0\) la escala y  \(\alpha=1/\xi>0\) la forma. 
	Es un caso particular de la distribución GEV cuando \(\xi>0\).
	
\[
f_{\alpha,\mu,\sigma}(x)=\frac{\alpha}{\sigma} \left(\frac{x-\mu}{\sigma}\right)^{-1-\alpha}  e^{-(\frac{x-\mu}{\sigma})^{-\alpha}}
\]

\begin{center}
%	\includegraphics[width=0.90\textwidth]{../Graph/LoiFrechet.pdf}
	
		\tikzset{
declare function={
	dFrechet(\x,\m,\sigma,\alpha) = (\alpha/\sigma) * (((\x-\m)/\sigma)^(-1-\alpha)) * exp(- ((\x-\m)/\sigma)^(-\alpha) );}
}  
\begin{tikzpicture}[mark=none,samples=100,smooth,domain=0:15]
\begin{axis}[ ytick=\empty,xtick pos=left,
	axis y line=middle,
	axis x line=bottom,
	height=.5*\linewidth,width=0.95*\linewidth, ticklabel style ={font=\footnotesize}, legend pos= north east,
	legend style={font=\tiny}]
	\addplot[BleuProfondIRA] 
	({\x},{dFrechet({\x},0,1,1)}) ;
	\addplot[FushiaIRA] 
	({\x},{dFrechet({\x},0,2,1)}) ;
	\addplot[VertIRA] 
	({\x},{dFrechet({\x},0,3,1)}) ;
	\addplot[OrangeProfondIRA] 
	({\x},{dFrechet({\x},0,2,2)}) ;
	%
	\legend{\(\mu\) =  0 \(\alpha=1\) \(\sigma = 1\),\(\mu\) =  0  \(\alpha=2\) \(\sigma =1\),\(\mu\) =  0  \(\alpha=3\) \(\sigma =1\),\(\mu\) = 0  \(\alpha=2\) \(\sigma = 2\)};
\end{axis}	 
\end{tikzpicture}
\end{center}


%\begin{center}
%%	\includegraphics[width=0.90\textwidth]{../Graph/FLoiFrechet.pdf}
%	{\shorthandoff{!;}	
%\tikzset{
%	declare function={
%		pFrechet(\x,\m,\sigma,\alpha) = exp( -((\x-\m)/\sigma)^(-\alpha));}
%}  
%\begin{tikzpicture}[mark=none,samples=100,smooth,domain=0:15]
%	\begin{axis}[ ytick=\empty,xtick pos=left,
%		axis y line=middle,
%		axis x line=bottom,
%		height=.5*\linewidth,width=0.95*\linewidth, ticklabel style ={font=\footnotesize}, legend pos= south east,
%		legend style={font=\tiny}]
%		\addplot[BleuProfondIRA] 
%		({\x},{pFrechet({\x},0,1,1)}) ;
%		\addplot[FushiaIRA] 
%		({\x},{pFrechet({\x},0,2,1)}) ;
%		\addplot[VertIRA] 
%		({\x},{pFrechet({\x},0,3,1)}) ;
%		\addplot[OrangeProfondIRA] 
%		({\x},{pFrechet({\x},0,2,2)}) ;
%		%
%		\legend{\(\mu\) =  0 \(\alpha=1\) \(\sigma = 1\),\(\mu\) =  0  \(\alpha=2\) \(\sigma =1\),\(\mu\) =  0  \(\alpha=3\) \(\sigma =1\),\(\mu\) = 0  \(\alpha=2\) \(\sigma = 2\)};
%	\end{axis}	 
%\end{tikzpicture}
%}
%\end{center}
\end{f}

\begin{f}[Relación entre GEV, Gumbel, Fréchet y Weibull]
	
	El parámetro de forma \(\xi\) determina el comportamiento de la cola de la distribución. 
	Las subfamilias definidas por \(\xi = 0\), \(\xi > 0\) y \(\xi < 0\) corresponden, respectivamente, a las familias de Gumbel, Fréchet y Weibull:
	\begin{itemize}
		\item      Gumbel o ley de valores extremos de tipo I
		\item     Fréchet o ley de valores extremos de tipo II, si \(\xi  = \alpha^{-1}\) con \(\alpha>0\),
		\item    Weibull invertida (\(\overline{F}\)) o ley de valores extremos de tipo III, si \(\xi =-\alpha^{-1} \), con \(\alpha>0\).
	\end{itemize}
\end{f}
\hrule 

\begin{f}[Teorema general de los valores extremos] 
	
	Sean \(X_1, \dots, X_n\) variables \(iid\), con función de distribución \(F_X\), y sea \(M_n =\max(X_1,\dots,X_n)\).
	
	La teoría proporciona la distribución exacta del máximo:
	\begin{align*}
		\mathcal{P}(M_n \leq z) = &\Pr(X_1 \leq z, \dots, X_n \leq z) \\
		= &\mathcal{P}(X_1 \leq z) \cdots \mathcal{P}(X_n \leq z) = (F_X(z))^n. 
	\end{align*}
	Si existe una secuencia de pares de números reales \((a_n, b_n)\) tal que \(a_n>0\) y \(\lim_{n \to \infty}\mathcal{P}\left(\frac{M_n-b_n}{a_n}\leq x\right) = F_X(x)\), donde \(F_X\) es una función de distribución no degenerada, entonces el límite de la función \(F_X\) pertenece a la familia de las leyes \(GEV\). 
	
\end{f}

\begin{f}[Densidad subexponencial]
	
	\textbf{Caso de las potencias} 
	
	Si \(\overline{F}_{X}(x)=\mathbb{P}(X > x)\sim c\ x^{-\alpha}\)
	cuando \(x \to \infty \) para un \(\alpha > 0 \) y una constante \(c > 0 \), entonces la ley de \(X\)
	es subexponencial.
	
	Si \(F_X\) es una función de distribución continua de esperanza \(\mathbb{E}[X]\) finita, se denomina índice de grandes riesgos a
	\[
	D_{F_X}(p)=\frac{1}{\mathbb{E}[X]}\int_{1-p}^{1} F_X^{-1}(t)dt,\, \, p\in [0,1]
	\]
	
	Esta distribución de exceso disminuye más lentamente que cualquier distribución exponencial.
	Es posible considerar esta estadística ~:
	\[
	T_n(p)=\frac{X_{(1:n)}+X_{(2:n)}+\ldots + X_{(np:n)}}{\sum_{1\leq i\leq n}(X_i)} \mbox{ donde } \frac{1}{n}\leq p\leq 1
	\]
	\(X_{(i:n)}\) denota el \(i\)-ésimo \(\max\) de los \(X_i\).
	
\end{f}
\begin{f}[Teorema de Pickands–Balkema–de Haan (ley de los excesos)]
	Sea \(X\) de distribución \(F_X\), y sea \(u\) un umbral elevado. Entonces, para una amplia clase de distribuciones \(F_X\), la distribución condicional de los excesos
	\[
	X_u := X - u \mid X > u
	\]
	es aproximable, para \(u\) suficientemente grande, por una ley de Pareto generalizada (GPD):
	
	\[
	\mathbb{P}(X - u \le y \mid X > u) \approx G_{\xi, \sigma, \mu=0}(y) :=  1 - \left(1+ \frac{\xi(x-\mu)}{\sigma}\right)^{-1/\xi} 
	\]
	
	\(\quad y \ge 0\). Dicho de otro modo, para \(u \to x_F := \sup\{x : F(x) < 1\}\),
	\[
	\sup_{0 \le y < x_F - u} \left| \mathbb{P}(X - u \le y \mid X > u) - G_{\xi,\sigma,\mu=0}(y) \right| \to 0.
	\]
	
	Este teorema justifica el uso de la \textit{ley de Pareto (generalizada)} para modelar los excesos por encima de un umbral, lo que constituye precisamente el marco de los contratos de reaseguro de \textit{exceso de pérdida} por riesgo, por evento o de acumulación anual. 
\end{f}
\hrule


\begin{f}[Datos sobre reaseguros]
	
	Dado que el reaseguro indemniza por acumulaciones de siniestros o siniestros extremos, suele recurrir a datos históricos que deben utilizarse con precaución~:
	\begin{itemize}
		\item la actualización de los datos (impacto de la inflación monetaria).
		\item la revalorización tiene en cuenta la evolución del riesgo:
		\begin{itemize}
			\item la evolución de las tasas de prima, las garantías y las condiciones de los contratos,
			\item la evolución de los costes de los siniestros (índice de costes de la construcción, índices de costes de reparación de automóviles, \ldots)
			\item la evolución del marco jurídico.
		\end{itemize}
		\item el ajuste de las estadísticas para tener en cuenta la evolución de la base de la cartera:
		\begin{itemize}
			\item perfil de las pólizas (número, capitales, \ldots),
			\item naturaleza de las garantías (evolución de las franquicias, de las exclusiones, \ldots)
		\end{itemize}
	\end{itemize}
	Tras estas correcciones, los datos se denominan «as if» (en economía, se utiliza la expresión «contrafactual»).

\end{f}


\begin{f}[La prima «Burning Cost»]
	
	\(X_{i}^{j}\) designa el \(i^{ésimo}\) siniestro  del año \(j\), como si estuviera actualizado, revalorizado y  ajustado,
	\(n^j\) el número de siniestros del año \(j\), \(c^{j}\) el coste para la aseguradora \(c\).
	La tasa pura según el método de \textbf{Burning Cost} viene dada por la fórmula~:
	\[
	BC_{puro}=\frac{1}{s}\sum_{j=1}^{n}\frac{c^{j}}{a_{j}}
	\] 
	El Burning Cost no es más que una media de los ratios \(S/P\) cruzados~: los siniestros a cargo del reasegurador sobre las primas recibidas por la cedente.
	La prima Burning Cost es, por tanto: \(P_{pura}=BC_{puro}\times a_{s+1}\).
	
	En el caso de un \(p\) XS \(f\)), \[
	c^{j}=\sum_{i=1}^{n^j}\max\left(  \left( X_{i}^{j}-f\right),p\right) 1\!\!1_{x^{j}\geq f}
	\]
	Si el seguro de vida calcula las tasas de prima en función del capital, el seguro de no vida utiliza como referencia el valor asegurado, mientras que el reaseguro toma como referencia el total de las primas de la cedente, denominado \textbf{base imponible}.
	Se denota por \(a_{j}\) la base de primas del año \(j\) y por \(a_{s+1}^{*}\) la base estimada para el año siguiente, donde \(s\) es el número de años de historial.
\end{f}



\begin{f}[El modelo de Poisson-Pareto]
	[Prima del XS o del XL]
	Sean \(p\) y \(f\) el alcance y la prioridad (franquicia) del XS, respectivamente, con el límite \(l=p+f\) (\(p\) XS \(f\)).
	
	La prima del XS viene dada por:
\[
\mathbb{E}\left[S_N\right]=\mathbb{E}\left[\sum_{i=1}^{N} Y_{i}\right]=\mathbb{E}[N]\times \mathbb{E}[Y]
\]
donde
\[
\mathbb{E}[Y]= l \mathbb{P}[X>l] - f \times \mathbb{P}[X\geq f] + \mathbb{E}[X\mid f\geq X\geq l]
\]


Si \(l=\infty\) y \(\alpha \neq 1\) :
\[
\mathbb{E}[S_N] = \lambda \frac{x_\mathrm{m}^{\alpha}}{\alpha -1}f^{1-\alpha} 
\]

si  \(l=\infty\) y \(\alpha = 1\) no hay solución.

Si \(l<\infty\) y \(\alpha \neq 1\) :
\[
\mathbb{E}[S_N] = \lambda \frac{x_\mathrm{m}^{\alpha}}{\alpha -1}\left( f^{1-\alpha} -l^{1-\alpha} \right)
\]

Si \(l<\infty\) y \(\alpha = 1\) :
\[
\mathbb{E}[S_N] = \lambda x_\mathrm{m} \ln \left(  \frac{1}{f}\right)
\]
\end{f}



\begin{f}
[El modelo Poisson-LogNormal]

Si \(X\) sigue una distribución \(\mathcal{L}\mathcal{N}orm(x_\mathrm{m}, \mu, \sigma)\), entonces \(X-x_\mathrm{m}\) sigue una distribución \(\mathcal{L}\mathcal{N}orm(\mu, \sigma)\)
De lo que se deduce:
\[
\mathbb{P}[X>f]=\mathbb{P}[X-x_\mathrm{m}>f-x_\mathrm{m}]=1-\Phi\left(\frac{\ln(f-x_\mathrm{m})-\mu}{\sigma}\right)
\]
\begin{align*}
	\mathbb{E}[X\mid X>f] \\ 
	= & \mathbb{E}\left[ X-x_\mathrm{m}\mid X-x_\mathrm{m}>f-x_\mathrm{m} \right]+x_\mathrm{m} \mathbb{P}[X>f]\\
	= & e ^{m+\sigma^2/2} \left[1-\Phi\left(\frac{\ln(f-x_\mathrm{m})-(\mu+\sigma^2)}{\sigma}\right) \right]\\
	& +x_\mathrm{m} \left( 1-\Phi\left(\frac{\ln(f-x_\mathrm{m})-\mu}{\sigma}\right) \right)
\end{align*}
Con franqueza y sin límites:
\begin{align*}
	\mathbb{E}[S_N] \\ 
	= & \lambda \left(\mathbb{E}\left[ X-x_\mathrm{m}\mid X-x_\mathrm{m}>f-x_\mathrm{m} \right]+x_\mathrm{m} \mathbb{P}[X>f]-f\mathbb{P}[X>f]\right)\\
	= & \lambda \left(  e ^{m+\sigma^2/2} \left[1-\Phi\left(\frac{\ln(f-x_\mathrm{m})-(\mu+\sigma^2)}{\sigma}\right) \right]\right)\\
	& +\lambda (x_\mathrm{m}-l) \left( 1-\Phi\left(\frac{\ln(f-x_\mathrm{m})-\mu}{\sigma}\right) \right)
\end{align*}

\end{f}